% Paper 3, §10 — Search-vs-RL under fixed compute
% Anchors: kb/notes/scaling/inference-time-compute-scaling.md §2.2-2.3
%          kb/notes/post-training/rlvr-and-grpo.md §2.2
%          kb/excerpts/{snell2024,dapo2025,deepseek-r1}.md

\section{Search-vs-RL under fixed compute}
\label{sec:search-vs-rl}

\Cref{sec:inference-compute-scaling} (the Snell anchor) settled one
slice of the triangle's central operational tradeoff: at fixed total
compute $C_\text{train} + C_\text{test}$, when does spending the
marginal FLOP on inference-time search beat spending it on
\emph{pretraining}? The empirical answer is conditional on difficulty
regime: on easy and intermediate prompts test-time compute substitutes
for $\sim\!14\times$ more pretraining FLOPs at matched accuracy
\citep[\S 6]{snell2024};\footnote{KB excerpt:
\texttt{kb/excerpts/snell2024.md\#sec-1-flops-matched}.} on the hardest
tail it does not. \Cref{sec:rlvr-grpo-dapo} introduced a third
allocation channel that did not exist when Snell's curves were drawn:
RL on verifiable rewards (RLVR) at GRPO/DAPO scale, which spends
FLOPs on \emph{post-training} the proposer rather than on either
pretraining the proposer or running it longer at inference. The
triangle's central operational question therefore generalises:
\textbf{at fixed total compute, when does the marginal FLOP go to
(a) more inference search/sampling, (b) more RLVR training, or (c)
more pretraining?} This section transcribes what is settled, isolates
what is anecdotal, and demarcates the open frontier.

\subsection{The three-channel allocation}
\label{sec:svr-three-channels}

Total compute spent to produce a deployed answer to a single prompt
decomposes into three additive components:
\begin{equation}
  C_\text{total}
    \;=\; \underbrace{C_\text{pre}}_{\text{pretraining}}
    \;+\; \underbrace{C_\text{RL}}_{\text{RLVR / GRPO post-training}}
    \;+\; \underbrace{C_\text{test}}_{\text{inference search/sampling}}.
  \label{eq:three-channel}
\end{equation}
The three terms are amortised differently. $C_\text{pre}$ is paid once
per base model and amortises across every prompt the model ever sees.
$C_\text{RL}$ is paid once per post-training run and amortises across
every prompt the post-trained checkpoint serves. $C_\text{test}$ is
paid \emph{per prompt at inference} and does not amortise. The
appropriate accounting for a budget question therefore depends on the
deployment volume: a one-shot research evaluation amortises
$C_\text{pre}$ and $C_\text{RL}$ over very few queries and is
inference-cheap; a productionised consumer endpoint amortises both over
$10^9$+ queries and is inference-bound.

Snell's FLOPs-matched experiment of \cref{eq:flops-matched} is the
two-channel restriction $C_\text{RL} = 0$:
\begin{equation}
  C_\text{small}^\text{train} + C^\text{test}
    \;\approx\; C_\text{large}^\text{train} + C_\text{large}^\text{1-pass}.
  \label{eq:svr-flops-matched-2chan}
\end{equation}
\citet[\S 6]{snell2024} compare a smaller base model running
compute-optimal test-time strategies against a $14\times$ larger base
model running a single inference pass at matched total FLOPs. The
match is between $C_\text{pre}$ on the right and
$C_\text{pre} + C_\text{test}$ on the left; the RL channel does not
appear because PaLM~2-S\* (Codey) is not RLVR-trained
\citep[\S 4]{snell2024}.\footnote{KB excerpt:
\texttt{kb/excerpts/snell2024.md\#sec-4-setup}.} What
\cref{sec:rlvr-grpo-dapo} added is a second axis of post-training
compute that was not in the original FLOPs-matched menu.

\begin{intuition}
The folk gloss of \cref{eq:three-channel} is ``where do the next
million GPU-hours go?'' The folk gloss is exactly right; the only
subtlety is that the answer depends on the deployment volume and the
prompt-difficulty distribution simultaneously. Returning to math: each
of $C_\text{pre}, C_\text{RL}, C_\text{test}$ is a knob in
\cref{eq:three-channel}, and the joint argmax over the three at fixed
$C_\text{total}$ subject to a target accuracy on a target prompt
distribution is the natural generalisation of Snell's
\cref{eq:snell-compute-optimal} from a $\theta$-only argmax at fixed
$C_\text{test}$ to a $(\theta_\text{pre}, \theta_\text{RL},
\theta_\text{test})$ argmax at fixed $C_\text{total}$.
\end{intuition}

\subsection{Snell's pretraining-vs-inference slice}
\label{sec:svr-pretraining-vs-inference}

The Snell two-channel result establishes the substitution direction
between $C_\text{pre}$ and $C_\text{test}$
\citep[\S 6]{snell2024}.\footnote{KB excerpt:
\texttt{kb/excerpts/snell2024.md\#sec-1-flops-matched}.} At matched
total FLOPs, the smaller-base-plus-test-time configuration (the
left-hand side of \cref{eq:svr-flops-matched-2chan}) outperforms the
$14\times$-larger-base configuration on easy and intermediate
difficulty bins, breaks even or under-performs on the hard bin, and
materially under-performs on the very-hard tail where
$\widehat{\mathrm{pass@}1}(q)$ on the small base is near zero. The
substitution ratio is therefore not a single number; it is a
\emph{difficulty-conditional curve}, and the $14\times$ headline is
the easy-bin reading.

Two consequences for \cref{eq:three-channel} follow directly. First,
on the easy/intermediate tail, marginal FLOPs are better spent on
$C_\text{test}$ than on $C_\text{pre}$. Second, on the very-hard tail,
the order reverses: marginal FLOPs are better spent on $C_\text{pre}$
than on $C_\text{test}$. The cross-over is benchmark-specific and
base-model-specific; \citet[\S 6]{snell2024} report it for the MATH
benchmark on PaLM~2-S\*, and the cross-over has not been
re-measured on a frontier-2026 reasoning-trained base model
(\cref{sec:snell-open}).

\subsection{The RL-vs-inference slice}
\label{sec:svr-rl-vs-inference}

The third channel of \cref{eq:three-channel} requires partial
empirical answers from the RLVR literature, since no Snell-style
FLOPs-matched study has been published for $C_\text{RL}$ vs.\
$C_\text{test}$ as of 2026-Q2. Two anchor data points exist.

\paragraph{R1's training-token cost.} \citet[\S 2.1]{deepseek-r1}%
\footnote{KB excerpt:
\texttt{kb/excerpts/deepseek-r1.md\#sec-2-1-hp}.} report
DeepSeek-R1-Zero training as $10{,}400$ GRPO steps on $G = 16$ rollouts
per prompt with maximum trajectory length $32{,}768$ tokens (raised to
$65{,}536$ tokens after step $8{,}200$). The training-time token
budget is therefore order $10^{4} \cdot 16 \cdot 3.3 \cdot 10^{4}
\approx 5 \cdot 10^{9}$ tokens of generation per epoch (rough), times
$1.6$ epochs \citep[\S 2.1]{deepseek-r1}. AIME~2024 pass@1 climbs from
$15.6\%$ on DeepSeek-V3-Base to $77.9\%$ on R1-Zero
\citep[\S 2.3]{deepseek-r1};\footnote{KB excerpt:
\texttt{kb/excerpts/deepseek-r1.md\#sec-2-3-aime}.} self-consistency
at inference adds a further $8.8\%$ to $86.7\%$ at $N = 64$ samples
\citep[\S 2.3]{deepseek-r1}. The $N = 64$ inference branch costs
$\sim 64$ forward-pass-equivalents per prompt; the GRPO training
branch costs $\sim 5\cdot 10^{9}$ tokens, amortised across every
deployed query. For a deployment that serves $\gtrsim 10^{8}$ queries
the per-query share of $C_\text{RL}$ is below the per-query
$C_\text{test}$ at $N = 64$; for a one-shot evaluation it is not.

\paragraph{DAPO's step-efficiency.}
\citet[\S 4]{dapo2025}%
\footnote{KB excerpt: \texttt{kb/excerpts/dapo2025.md\#sec-4-headline}.}
report that DAPO reaches $50$ points on AIME~2024 with the
Qwen2.5-32B-base in roughly half the GRPO steps the DeepSeek-R1-Zero
recipe would require to reach a comparable score. The four targeted
modifications --- clip-higher, dynamic sampling, token-level loss,
overlong reward shaping --- compose multiplicatively on the
step-efficiency dimension \citep[\S 3]{dapo2025}. From the perspective
of \cref{eq:three-channel}, DAPO halves $C_\text{RL}$ at matched
post-training accuracy. The substitution direction is therefore: a
better RL recipe makes the RL channel cheaper without changing
$C_\text{pre}$ or $C_\text{test}$, and the saved $C_\text{RL}$ is
re-allocatable to either of the other two channels.

\paragraph{The composite R1 + self-consistency anecdote.} The R1
empirics combine the second and third channels. Starting from the
DeepSeek-V3-Base model (the $C_\text{pre}$ channel, fixed), GRPO with
verifiable rewards lifts AIME~2024 from $15.6\%$ to $77.9\%$ pass@1
(the $C_\text{RL}$ channel, $\sim\!10^{4}$ GRPO steps); self-consistency
at $N = 64$ then adds $8.8\%$ to $86.7\%$ (the $C_\text{test}$ channel,
$N$-fold inference). Reading the curve as a sequence of marginal
returns: the first units of $C_\text{RL}$ deliver the largest accuracy
gain ($+62.3$ percentage points from $15.6\%$ to $77.9\%$); the first
units of $C_\text{test}$ on top of the trained model deliver an
incremental gain of $+8.8$ points at the cost of $64\times$ inference
flops per query. The marginal-FLOP ranking is anecdotal and
benchmark-specific, but the qualitative pattern --- RL channel buys
the bulk of the gain on math, inference channel buys the long tail of
sample-efficient sharpening --- recurs across the 2025-2026
reasoning-trained-model literature \citep[\S 2.3]{deepseek-r1}.

\subsection{The identity that ties the channels}
\label{sec:svr-identity}

The two-channel \cref{eq:svr-flops-matched-2chan} generalises in
the obvious way to three channels. Holding total compute fixed,
\begin{equation}
  C_\text{small}^\text{train} + C_\text{RL}^\text{small} + C^\text{test}
    \;\approx\; C_\text{large}^\text{train} + C_\text{large}^\text{1-pass},
  \label{eq:svr-three-channel-identity}
\end{equation}
where the left side is a smaller base model post-trained with RLVR
plus inference-time search, and the right side is the
$14\times$-larger un-RL'd base model running a single inference pass.
\Cref{eq:svr-three-channel-identity} is a budget bookkeeping
identity, not a scaling law: it commits to no functional form for
how accuracy depends on the three components on either side. The
empirical question is whether (and on which prompt-difficulty bins)
the left-hand configuration matches the right-hand accuracy. As of
2026-Q2 the three-channel version of the FLOPs-matched experiment
has not been published; the two-channel restriction is Snell's, and
the third channel's empirical content is the anecdotes of
\cref{sec:svr-rl-vs-inference}.

\paragraph{What a closed-form would look like.} A Chinchilla-style
joint frontier in the three channels would be a parametric loss
\begin{equation}
  L(N, D, C_\text{RL}, C_\text{test})
    \;=\; E
      \;+\; \frac{A}{N^{\alpha}}
      \;+\; \frac{B}{D^{\beta}}
      \;+\; \Phi(C_\text{RL})
      \;+\; \Psi(C_\text{test}, q),
  \label{eq:svr-hypothetical-frontier}
\end{equation}
generalising \citet[\S 3, Eq.~2]{hoffmann2022-chinchilla}'s
$L(N, D) = E + A/N^\alpha + B/D^\beta$ with two additional terms
$\Phi$ and $\Psi$ capturing RL-channel and inference-channel returns.
The functional forms of $\Phi$ and $\Psi$ are unknown:
\citet{s1-2025}'s budget-forcing experiment is suggestive of
roughly logarithmic returns to thinking-token count for $\Psi$ on a
small reasoning-trained model (\cref{sec:snell-open}), but no
parametric fit with the structural status of Chinchilla's law has
been attested for either $\Phi$ or $\Psi$. The cross-derivative
$\partial^2 L / \partial C_\text{RL}\, \partial C_\text{test}$ ---
whether the RL and inference channels substitute or complement at
the margin --- is the empirical content
\cref{sec:svr-rl-vs-inference}'s anecdotes hint at, but no
closed-form has been fitted.

\subsection{Where the curves cross is benchmark-specific}
\label{sec:svr-benchmark-specific}

The three-channel allocation question is not solvable in the abstract.
The cross-over points among $\{C_\text{pre}, C_\text{RL},
C_\text{test}\}$ depend on at least four contingencies, each of which
shifts the curves measurably:

\begin{enumerate}
\item \textbf{Prompt-difficulty distribution.} Easy prompts favour
$C_\text{test}$ (Snell's revision regime); intermediate prompts favour
$C_\text{test}$ in search form (best-of-$N$-with-PRM, beam, MCTS); the
hardest tail favours $C_\text{pre}$ \citep[\S 6]{snell2024}. RLVR
training shifts the difficulty distribution itself: an
RL-trained model's difficulty bins are not the base model's bins
\citep[\S 2.3]{deepseek-r1}, which means the prompt-difficulty contingency
interacts with $C_\text{RL}$ rather than being orthogonal to it.

\item \textbf{Verifier availability.} The whole RLVR channel exists
because a programmatic verifier can be applied at training time
\citep[\S 1.1]{shao2024}. Domains without a verifier (open-ended
generation, agentic tasks without a programmatic check) cannot spend
the $C_\text{RL}$ channel directly, so the three-channel allocation
collapses to Snell's two-channel \cref{eq:svr-flops-matched-2chan}.
The frontier extension --- LLM-as-judge, generative PRMs as soft
verifiers (\cref{sec:prm-after-r1}) --- is exactly the attempt to
make $C_\text{RL}$ available outside the verifiable domain at the
cost of soft-verifier reward-hacking risk.

\item \textbf{Base-model checkpoint.} The substitution ratios are
measured against a specific base model. PaLM~2-S\* on MATH gives
$14\times$ at the easy bin \citep[\S 6]{snell2024}; the same
experiment on Qwen2.5-32B-base, on Llama 4, or on the still-classified
o1-base would produce different ratios. \citet[\S 6]{snell2024}
explicitly note their result depends on ``the specific conditions on
the pretraining and inference workload.''

\item \textbf{Deployment-volume amortisation.} The per-query share of
$C_\text{pre}$ and $C_\text{RL}$ scales as $1/(\text{queries served})$
while $C_\text{test}$ does not amortise. For a research evaluation,
$C_\text{test}$ dominates the per-query budget; for a frontier
production endpoint, $C_\text{pre}$ and $C_\text{RL}$ are
nearly-free-at-the-margin and $C_\text{test}$ is the binding
constraint.
\end{enumerate}

The four contingencies multiply, not add. A statement of the form
``the optimal split between RL training and inference search is X:Y''
is therefore meaningful only when scoped to all four: a specific
prompt-difficulty distribution, a specific verifier availability, a
specific base model, and a specific deployment volume. The 2025-2026
literature reports such scoped statements --- the R1 anecdote of
\cref{sec:svr-rl-vs-inference} is one --- but does not extract a
universal exchange rate.

\begin{contradiction}
\textbf{No closed-form three-channel tradeoff curve is attested.}
\citet[\S 6]{snell2024} establish the pretraining-vs-inference
substitution at $14\times$ on the easy/intermediate MATH bin under
PaLM~2-S\*. \citet[\S 2.3]{deepseek-r1} establish that GRPO on
DeepSeek-V3-Base converts $\sim\!10^{4}$ training steps into $+62.3$
points of AIME~2024 pass@1 with self-consistency adding a further
$+8.8$ at $N = 64$. \citet[\S 4]{dapo2025} establish that DAPO halves
$C_\text{RL}$ at matched post-training accuracy. The three results
are anchor anecdotes, each on a different base model and benchmark.
No published study has run a three-channel FLOPs-matched experiment
holding $C_\text{total} = C_\text{pre} + C_\text{RL} + C_\text{test}$
fixed and varying the allocation, and no parametric fit
to \cref{eq:svr-hypothetical-frontier} with the structural status of
Chinchilla's $L(N, D) = E + A/N^\alpha + B/D^\beta$ exists. Whether a
joint $(C_\text{train}, C_\text{RL}, C_\text{test})$ frontier with a
clean closed form exists at all is open; the contingencies of
\cref{sec:svr-benchmark-specific} suggest the answer may be ``only
within a fixed prompt-difficulty distribution, verifier setting, base
model, and deployment regime,'' which is a much weaker statement than
the universal Chinchilla-style law it would generalise. Re-visited in
\cref{sec:contradictions}.
\end{contradiction}

\subsection{Operating-point heuristics from the anchor anecdotes}
\label{sec:svr-heuristics}

Although no closed-form tradeoff is attested, the anchor results of
\cref{sec:svr-pretraining-vs-inference,sec:svr-rl-vs-inference}
support three operating-point heuristics that are stable across the
2025-2026 reasoning literature:

\begin{itemize}
\item \textbf{If the verifier is available, spend on $C_\text{RL}$
first.} The R1 anecdote shows the first units of GRPO buy the largest
accuracy gain on math benchmarks
\citep[\S 2.3]{deepseek-r1}.\footnote{KB excerpt:
\texttt{kb/excerpts/deepseek-r1.md\#sec-2-3-aime}.} The marginal
return on $C_\text{RL}$ exceeds the marginal return on either
$C_\text{test}$ (small inference samples on an un-RL'd base) or
$C_\text{pre}$ (Snell's $14\times$ pretraining substitution on an
un-RL'd base) until the policy saturates the verifier-detectable
gain.
\item \textbf{Once the policy is RL-trained, $C_\text{test}$ delivers
incremental sharpening.} Self-consistency on R1-Zero adds $8.8$
points at $N = 64$ \citep[\S 2.3]{deepseek-r1}; the same algorithm
on the un-RL'd base would not approach this number. The
inference-channel return depends on the base distribution it operates
on (\cref{sec:self-consistency,sec:inference-compute-scaling}),
and the RL-trained distribution is sharper.
\item \textbf{The very-hard tail wants $C_\text{pre}$.} For prompts
the base model essentially never solves
($\widehat{\mathrm{pass@}1}(q) \to 0$),
\citet[\S 6]{snell2024} report no inference strategy recovers the
gap. RLVR cannot bootstrap on a $G$-group of trajectories that all
score $R_i = 0$ either, because the group-normalised advantage of
\cref{eq:grpo-advantage} is zero by construction
\citep[\S 4.1]{shao2024}; this is the failure mode DAPO's dynamic
sampling addresses by discarding all-zero groups
\citep[\S 3]{dapo2025}\footnote{KB excerpt:
\texttt{kb/excerpts/dapo2025.md\#sec-3-dynamic-sampling}.} but cannot
solve at the level of the underlying capability gap. The very-hard
tail therefore requires expanding the base distribution's support,
which is the $C_\text{pre}$ channel.
\end{itemize}

The three heuristics are scoped: they hold for verifiable-reward
domains, for the prompt-difficulty bins where a verifier discriminates,
and on the timescales where amortisation favours one-time training
costs over per-query inference costs. None of them is a substitute for
a closed-form tradeoff curve; they are operating-point recipes that
the 2026 frontier deploys.

\begin{forumsignal}
The shorthand ``RL is more compute-efficient than search'' is
sometimes cited as a derivable consequence of the R1 anecdote. It is
not. The R1 anecdote shows that on \emph{a specific base model and
benchmark, with a specific verifier and a specific deployment-volume
amortisation}, the first units of $C_\text{RL}$ deliver more accuracy
per FLOP than the first units of $C_\text{test}$ on top of the un-RL'd
base. The shorthand drops every contingency. Cited or repeated
without scope, it should be treated as a discovery signal that RL
buys a lot on math under verifiers, not as a numerical exchange rate
transferable to non-verifiable domains, frontier base models, or
research-scale deployments. \citet[\S 6]{snell2024} make the same
point about the $14\times$ substitution ratio.
\end{forumsignal}

\subsection{Cross-paper and forward links}
\label{sec:svr-forward}

The three-channel allocation question is the operational form of the
triangle's central tradeoff. \Cref{sec:cot-faithfulness} (CoT
faithfulness) is structurally adjacent: it asks what causal work the
$C_\text{test}$ channel's reasoning trace actually does in producing
the answer, which is upstream of any quantitative split between
$C_\text{RL}$ and $C_\text{test}$ --- if the trace is post-hoc
rationalisation \citep[KB note]{arcuschin2025-cot-wild},
\deepencite{arcuschin2025-cot-wild §3, citation-pending}
the inference-channel return is upper-bounded by what answer-only
sampling could achieve. \Cref{sec:r1-family} (the R1 family,
forthcoming) makes \cref{sec:svr-rl-vs-inference}'s composite
anecdote concrete by walking the four-stage R1 protocol; the
empirics of $C_\text{RL}$ in this section are taken from that
protocol's first RL stage. \Cref{sec:frontier-2026} (frontier
snapshot, forthcoming) catalogues the production combinations of the
three channels deployed by R1, o1/o3, Gemini~2.5-thinking,
Claude~3.7-extended-thinking, Qwen3-thinking, and QwQ; the snapshot
is the empirical census the closed-form
\cref{eq:svr-hypothetical-frontier} would have to fit.
\Cref{sec:contradictions} (contradictions live and open frontiers)
collects the open-frontier reading of this section's central
contradiction --- whether a joint
$(C_\text{train}, C_\text{RL}, C_\text{test})$ scaling law exists ---
alongside the other five \texttt{[CONTRADICTION]} markers carried
forward from \cref{sec:cot-phenomenon,sec:tree-search,sec:rlvr-grpo-dapo,sec:r1-family,sec:cot-faithfulness,sec:prm-after-r1}.

The structural picture: \cref{eq:three-channel} is the budget
identity; \cref{eq:svr-three-channel-identity} is its FLOPs-matched
form; \cref{eq:svr-hypothetical-frontier} is the parametric law that
would close the loop and currently does not exist. The triangle's
three legs are individually well-characterised
(\cref{sec:inference-compute-scaling,sec:rlvr-grpo-dapo,sec:r1-family}),
but their joint frontier is the open research direction this section
demarcates. The closing thesis of \cref{sec:contradictions} returns
to this frontier as the most consequential of the paper's six live
contradictions for the next model generation's compute-allocation
decisions.
